\chapter{Conclusiones y trabajo futuro}
\label{cap:conclusiones}

Este trabajo ha abordado el problema de verificar la comprensión real que un
estudiante tiene del código que entrega, un reto acentuado por la facilidad con
la que hoy puede generarse código mediante herramientas de inteligencia
artificial. Como respuesta, se ha diseñado, implementado y evaluado una
herramienta que genera de forma automática exámenes escritos y personalizados
sobre el código de cada alumno. Este capítulo recoge el grado de cumplimiento
de los objetivos, las limitaciones detectadas y las líneas de trabajo futuro.


\section{Cumplimiento de los objetivos}
\label{sec:cumplimiento-objetivos}

El \textbf{objetivo general} ---diseñar, implementar y evaluar una herramienta
que genere un examen personalizado a partir del código entregado--- se ha
alcanzado: el sistema recorre una carpeta de código, extrae sus unidades,
genera preguntas de distintos tipos revisadas por un sistema multi-agente y
produce el examen final. En cuanto a los objetivos específicos:

\begin{itemize}
  \item \textbf{Análisis de código (objetivo 1).} Se ha implementado un módulo
        basado en \emph{tree-sitter} que extrae funciones, métodos y clases de
        una carpeta completa y descarta las unidades triviales. Su
        arquitectura, organizada como un registro de analizadores por lenguaje,
        demostró ser extensible: se añadió soporte para Java sin tocar
        el resto del sistema. A raíz de la validación con el docente se
        incorporó, además, la capacidad de descartar el código de plantilla, de
        modo que solo se evalúa el código que el alumno ha escrito.
  \item \textbf{Generación multi-agente (objetivo 2).} Se ha construido un
        sistema \emph{generador-revisor} que produce preguntas de opción
        múltiple, de traza y de respuesta abierta, adaptadas al código y al
        nivel académico configurado.
  \item \textbf{Estudio de ablación (objetivo 3).} Se ha comparado la calidad
        de las preguntas con y sin el agente revisor. El revisor mejora la
        calidad del conjunto entregado ---del 82\,\% al 96\,\% de preguntas
        correctas---, y es especialmente eficaz en Python.
  \item \textbf{Interfaz y salida (objetivo 4).} Se ha desarrollado una
        interfaz web que permite aportar el código, configurar el contexto
        académico y obtener el examen, con salida en PDF y también en Word
        editable.
  \item \textbf{Validación (objetivo 5).} La herramienta se ha probado sobre
        código real de prácticas de asignaturas de informática, en dos
        lenguajes, y sus preguntas se han sometido al criterio de un docente en
        activo.
\end{itemize}


\section{Aportaciones y hallazgos}
\label{sec:aportaciones}

Además de la propia herramienta, el trabajo deja algunas conclusiones
interesantes. El estudio de ablación aporta evidencia de que incorporar un segundo
agente revisor mejora la calidad de las preguntas generadas. Al mismo tiempo,
tanto la evaluación automática ---en la que un modelo juez se dejaba errores
técnicos sutiles--- como la validación con el docente ---que detectó problemas
de contexto que el sistema no podía ver por sí solo, como el uso de código de
plantilla--- coinciden en lo mismo: la herramienta rinde mejor como
\textbf{apoyo al docente} que como sustituto, dejando en sus manos la decisión
final sobre las preguntas.


\section{Limitaciones}
\label{sec:limitaciones}

El trabajo presenta varias limitaciones. La calidad de las preguntas es
\textbf{menor en Java que en Python}, lo que indica que tanto la generación
como la revisión están mejor ajustadas a este último. La evaluación automática
de la corrección mediante un modelo juez \textbf{no es del todo fiable} en los
errores más sutiles. La validación de utilidad es \textbf{todavía limitada}, al basarse en el
criterio de un solo docente y sobre una muestra reducida. Por último, el sistema genera preguntas sobre \textbf{unidades de
código aisladas} y no evalúa la comprensión de las relaciones entre ellas.


\section{Trabajo futuro}
\label{sec:trabajo-futuro}

De las limitaciones anteriores y del propio desarrollo surgen varias líneas de
trabajo futuro:

\begin{itemize}
  \item \textbf{Comprensión global del proyecto.} Extender el sistema para
        generar preguntas sobre las relaciones entre unidades ---llamadas entre
        funciones, colaboración entre clases--- y no solo sobre unidades
        aisladas.
  \item \textbf{Soporte de más lenguajes.} Ampliar el registro de analizadores
        con nuevos lenguajes y, sobre todo, mejorar la calidad de la generación
        y la revisión en lenguajes distintos de Python.
  \item \textbf{Entrada desde repositorios remotos.} Permitir tomar el código
        directamente de una URL ---por ejemplo, un repositorio de GitHub---,
        además de desde una carpeta local.
  \item \textbf{Ajuste del formato del examen.} Permitir configurar cuánto
        código se muestra en cada pregunta ---para no incluir más del necesario
        en un examen tradicional--- y afinar las preguntas de traza para que se
        centren en los aspectos relevantes del código y no en detalles menores.
  \item \textbf{Estudio sistemático de los \emph{prompts}.} Analizar de forma
        controlada cómo afectan distintas formulaciones de las instrucciones a
        la calidad de las preguntas.
  \item \textbf{Definición explícita de los niveles de dificultad.} Concretar
        qué implica cada nivel (básico, intermedio, avanzado) mediante una
        rúbrica o ejemplos calibrados, en lugar de delegar por completo su
        interpretación en el modelo, para que la dificultad sea más consistente
        y predecible.
  \item \textbf{Validación a gran escala.} Someter la herramienta al criterio
        de más docentes y sobre más código, para respaldar los resultados con
        una muestra mayor.
\end{itemize}


\section{Valoración final}
\label{sec:valoracion-final}

En conjunto, el trabajo ha cumplido sus objetivos: se ha construido una
herramienta funcional que genera exámenes personalizados sobre el código de
cada estudiante y se ha evaluado su comportamiento. Los resultados muestran que
la herramienta es viable y útil como apoyo al docente y, a la vez, dejan claro
por dónde puede seguir mejorando.
